\section{Introduction}
\label{sec:introduction}

Large language models (LLMs) are increasingly deployed as interactive services
where many users submit prompts concurrently and expect low-latency responses.
Because decoding is autoregressive, a request's total service time is
proportional to the number of tokens it generates, and that number is not
known when the request arrives. Serving systems that admit requests in
first-come-first-served (FCFS) order therefore have no way to distinguish a
prompt that will generate a handful of tokens from one that will generate
several hundred. When a long-running request is admitted into a batch ahead of
several short ones, the short requests are forced to wait behind it for the
duration of the long request's decoding, a phenomenon commonly referred to as
head-of-line (HOL) blocking. HOL blocking directly inflates tail latency and
degrades the perceived responsiveness of the service, even when aggregate GPU
throughput remains high.

A recent line of work addresses this problem by learning to predict, rather
than assume, a request's likely output length ahead of time, and using that
prediction to reorder the admission queue so that requests expected to finish
quickly are not stuck behind requests expected to run long. One such approach
trains a compact predictor -- a fine-tuned OPT-125M model -- with a
learning-to-rank (LTR) objective to produce a relative ranking score for each
incoming prompt, and schedules requests in order of that score. This is
effective when the predictor's training distribution matches the requests it
sees in production, but the score is a single learned signal: if the predictor
is uncertain or the incoming prompts drift from its training distribution, the
resulting ranking can degrade with no fallback signal to correct it.

This project asks whether a second, inexpensive, and non-learned signal can
make the ranking more robust without discarding the learned predictor. The
tokenized length of a prompt is available immediately, requires no additional
model inference, and does not depend on the predictor's calibration. We
combine it with the LTR predictor's score through a weighted composite formula
and integrate the resulting scheduler into vLLM, a widely used open-source LLM
serving engine, as a drop-in replacement for the admission-ordering step used
by the original LTR scheduler. We refer to the resulting scheduler as
\emph{LTR + Prompt Length} throughout this paper.

The contributions of this work are:
\begin{itemize}
    \item A composite ranking score that augments a learned LTR predictor
    score with tokenized prompt length as a lightweight, always-available
    secondary signal, requiring no additional model inference at scheduling
    time.
    \item An implementation of the composite scheduler inside vLLM's offline
    serving pipeline, evaluated alongside FCFS, a classification-based
    baseline, and the original LTR scheduler under identical hardware and
    workload conditions.
    \item An empirical evaluation on two 1{,}000-request subsets of
    LMSYS-Chat-1M -- one in-distribution and one constructed to induce
    distribution shift -- covering throughput, tail-latency indicators of HOL
    blocking, time-to-first-token, GPU resource utilization, and ranking
    quality via Kendall's tau.
    \item A transparent discussion of the correction applied to an initial
    sign error in the composite formula, the software-version deviations
    required to run the pipeline on current hardware, and the limitations of
    the evaluation methodology.
\end{itemize}

The remainder of this paper is organized as follows. Section~\ref{sec:background}
reviews the serving architecture and scheduling concepts this work builds on.
Section~\ref{sec:related_work} situates our approach relative to existing LLM
serving and scheduling research. Section~\ref{sec:methodology} describes the
composite scoring formula, the scheduling policies compared, and the vLLM
integration used to run them. Section~\ref{sec:evaluation} presents the
experimental results.
Section~\ref{sec:discussion} interprets the findings and discusses limitations
and future work, and Section~\ref{sec:conclusion} concludes.
